\documentclass[12pt,letterpaper]{article}
\usepackage[T1]{fontenc}
\usepackage{anysize}
\usepackage{tikz}
\usepackage{amssymb}
\usepackage{comment}
\usepackage{graphicx}

\setlength{\parindent}{0cm}
\setlength{\parskip}{1em}

\newcommand{\contradiction}{%
\begin{tikzpicture}[rotate=45,x=0.5ex,y=0.5ex]
\draw[line width=.2ex] (0,2) -- (3,2) (0,1) -- (3,1) (1,3) -- (1,0) (2,3) -- (2,0);
\end{tikzpicture}
}

\marginsize{2cm}{2cm}{1cm}{1cm}

\begin{document}

\begin{titlepage}
    \vspace*{4cm}
    {\huge \center
        CS 325 Project 1 Report\\[1cm]
    }
    \center
    {\large 
        Group 3

        Date: \today 

    \textbf{Contributors:}
	Cezary Wojcik,
    Sean McGlothlin,
    Matthew Eilertson
	}

\end{titlepage}

\section{Algorithm Pseudocode}
\subsection{Algorithm 1}
\begin{verbatim}
function alg1(equations):
  for each line in equations:
    line is visible.
  for distance from 2 to length(equations):
    for j from 0 to (length(equations) - distance):
      k = j + distance
      for i from j + 1 to k:
        xjk, yjk = intersection(equations[j], equations[k])
        if equations[j].value(at point xjk) > equations[i].value(at point xjk):
          equations[i] is not visible.
\end{verbatim}

Asymptotic Running Time: $\Theta(n^3)$

\subsection{Algorithm 2}
\begin{verbatim}
function alg2(equations):
  for each line in equations:
    line is visible.
  for distance from 2 to length(equations):
    for j from 0 to (length(equations) - distance):
      k = j + distance
      for i from j + 1 to k:
        if equations[i] already visible:
          continue
        xjk, yjk = intersection(equations[j], equations[k])
        if equations[j].value(at point xjk) > equations[i].value(at point xjk):
          equations[i] is not visible.
\end{verbatim}

Asymptotic Running Time: $O(n^3)$

\newpage

\subsection{Algorithm 3}
\begin{verbatim}
function alg3(equations):
	for each line in equations:
	 line is visible
    for each line in equations starting at 3:
        // starting at the last one -> Yjk should be Yi-1
        for each visible(true) line starting at the last one:
            // check the intersection of visible line(k-1) and visible line(k). 
            // Compare y(x(k-1,k)) value to y of new line, yi(x(k-1,k)). 
            // If y value of intersection of k-1,k < the y value of yi at the 
            // intersection k-1,k then mark k as not visible.
            xkK, ykK = intersection(equations[k-1], equations[k])
            // if the intersection is below yi, then yjk is not visible,
            // and we should iterate backwards to find one that is.
            if ykK < equations[i].value(at point xkK):
                mark line k as not visible.
            else break
\end{verbatim}

Asymptotic Running Time: $\Theta(n)$

\section{Proof of Claim 2}

Suppose $Y_t = \{ y_{j_1}, ..., y_{j_k} \}$ is the visible subset of $Y_{i-1} = \{ y_1, ..., y_{i-1} \}$ where $t \leq i - 1$. Let $y_{j_k} \in Y_t$ be the last line in $Y_t$ such that $y_{j_k}(x_{j_k, j_{k-1}}) \geq y_i(x_{j_k, j_{k-1}})$ where $(x_{j_k, j_{k-1}}, y_{j_k}(x_{j_k, j_{k-1}}))$ is the intersection of lines $y_{j_k}$ and $y_{j_{k-i}}$. 

Let $Y_k = \{ y_{j_1}, ..., y_{j_k} \}$ where $k \leq t$ and $Y_k \subseteq Y_t$ and let $Y_i = \{ y_1, ..., y_i \}$. 

Let $Y_{k+i} = Y_k \cup \{ y_i \}$.

Suppose that $Y_{k+i}$ is \textit{not} the visible subset of $Y_i$. Then, $\exists y_m \in Y_{k+i}$ for which $y_m$ is not visible in $Y_i$. 

We know that $y_1$ and $y_i$ are visible in $Y_i$ because they have the lowest and highest slopes in $Y_i$. Therefore, since $y_m$ is not visible, $y_m$ can't be $y_i$, so $y_m \in Y_t$.

Since $y_m$ is not visible we know that $\forall y_p \in Y_i, \not\exists x$ such that $y_m(x) \geq y_p(x)$. However, since $Y_t$ is the visible subset of $Y_{i-1}$, we know that $y_m$ is visible in $Y_{i-1}$. Therefore, $\forall y_p \in Y_{i-1}, \exists x$ such that $y_m(x) \geq y_p(x)$. Since $Y_{i-1} \subset Y_{i}$, $y_i$ must cover up $y_m$ in $Y_i$ where it was visible in $Y_{i-1}$. Thus, $y_i(x_{m, m-1}) > y_m(x_{m, m-1})$ where $(x_{m, m-1}, y_m(x_{m, m-1}))$ is the intersection of lines $y_m$ and $y_{m-1}$.

We also know that $y_{j_k}$ is also visible due to its definition, so $m < k$. However, we defined $y_{j_k}$ to be the \textit{last} line in $Y_t$ such that $y_{j_k}(x_{j_k, j_{k-1}}) \geq y_i(x_{j_k, j_{k-1}})$ where $(x_{j_k, j_{k-1}}, y_{j_k}(x_{j_k, j_{k-1}}))$ is the intersection of lines $y_{j_k}$ and $y_{j_{k-i}}$.

Contradiction! \contradiction

Therefore, by contradiction, if $Y_t = \{ y_{j_1}, ..., y_{j_k} \}$ is the visible subset of $\{ y_1, ..., y_{i-1} \}$ where $t \leq i - 1$, then $\{ y_{j_1}, ... y_{j_k}, y_i \}$ is the visible subset of $\{ y_1, ..., y_i \}$ where $y_{j_k} \in Y_t$ is the last line in $Y_t$ such that $y_{j_k}(x_{j_k, j_{k-1}}) > y_i(x_{j_k, j_{k-1}})$ where $(x_{j_k, j_{k-1}}, y_{j_k}(x_{j_k, j_{k-1}}))$ is the intersection of lines $y_{j_k}$ and $y_{j_{k-i}}$.

$\blacksquare$

\section{Proof of Algorithm 3}

For a set of lines, $Y_n = \{ y_1, ..., y_n \}$, we can determine the subset of visible lines using algorithm 3.

Base Case: If we only have 1 or 2 lines in the set, they are trivially visible.

Inductive Hypothesis: The set of lines $\{ y_{j_1}, ..., y_{j_t} \}$ is the visible subset of $Y_i = \{ y_1, ..., y_i \}$ where  $t \leq i - 1$.

Inductive Step: Consider the set of lines $Y_{i+1} = \{ y_1, ..., y_{i+1} \}$. Given the Inductive Hypothesis, we can determine the visible subset of $Y_{i+1}$ using Claim 2.

Therefore, by induction, algorithm 3 can correctly determine the subset of visible lines of a given set of lines, $Y_n = \{ y_1, ..., y_n \}$.


$\blacksquare$

\section{Benchmarking Results and Analysis}

\includegraphics[scale=.70]{graph.png}

In the image above, about all we can tell is that the performance of algorithm 1 is far worse than the performance of algorithms 2 and 3. The graph is not very useful when in a linear scale.

\includegraphics[scale=.70]{loglog.png}

Once we applied a logarithmic scale to both axes on the graph, the results are far more easy to see. 

According to the trendline, if we had tried to run a benchmark for algorithm 1 for 9,000 random equations, it would have taken almost 27 days to complete, whereas algorithm 3 was able to accomplish it in under half a second.

If each algorithm were to run for an hour, then:

\begin{itemize}
	\item algorithm 1 could process around 1050 equations
	\item algorithm 2 could process around 9566 equations
	\item algorithm 3 could process around 77,088,900 equations
\end{itemize}

The slope of the regression line in the log-log plot is:

\begin{itemize}
	\item 3.0052 for algorithm 1
	\item 2.2494 for algorithm 2
	\item 0.9862 for algorithm 3
\end{itemize}

3.0052 is pretty close to the 3 in $\Theta(n^3)$, the asymptotic runtime for algorithm 1. 2.2494 is below 3, so it fits $O(n^3)$, the asymptotic runtime for algorithm 2. Finally 0.9862 is pretty close to 1 in $\Theta(n^1)$, the asymptotic runtime for algorithm 3. All in all, the benchmarking results match our asymptotic runtime analysis.

\end{document}
